\documentclass[12pt]{article}

% Standard math and formatting packages
\usepackage[utf8]{inputenc}
\usepackage[T1]{fontenc}
\usepackage{amsmath,amssymb,amsthm}
\usepackage{mathrsfs}
\usepackage{geometry}
\usepackage{hyperref}
\usepackage{url}
\usepackage{cleveref}
\usepackage{enumitem}
\usepackage{booktabs}
\usepackage{tabularx}

\geometry{margin=1in}

% Theorem environments
\newtheorem{theorem}{Theorem}[section]
\newtheorem{proposition}[theorem]{Proposition}
\newtheorem{corollary}[theorem]{Corollary}
\newtheorem{lemma}[theorem]{Lemma}
\newtheorem{claim}{Claim}

\theoremstyle{definition}
\newtheorem{definition}[theorem]{Definition}
\newtheorem{axiom}{Axiom}

\theoremstyle{remark}
\newtheorem{remark}[theorem]{Remark}

% Custom commands
\newcommand{\calX}{\mathcal{X}}
\newcommand{\calM}{\mathcal{M}}
\newcommand{\calG}{\mathcal{G}}
\newcommand{\Sigmabar}{\bar{\Sigma}}
\newcommand{\fg}{\mathfrak{g}}

\begin{document}

\begin{titlepage}
\centering
\vspace*{2cm}

{\LARGE\bfseries Foundations of Invariant Evaluation:\\
A Conservative Framework for\\
Representation-Independent Reasoning\\[10pt]
\large Grounding, Canonicalization, and Invariance as\\
Design Principles for Trustworthy Systems}

\vspace{1.5cm}

{\Large Andrew H. Bond, Senior Member, IEEE}

\vspace{0.5cm}

{\large
Department of Computer Engineering\\
San Jos\'e State University\\
San Jos\'e, CA 95192, USA\\[0.3cm]
\texttt{andrew.bond@sjsu.edu}\\[0.3cm]
ORCID: 0009-0003-2599-6158
}

\vspace{1cm}

{\large Spring 2026}

\vspace{2cm}

\begin{abstract}
We develop a framework for evaluation systems---whether ethical, safety-critical, or
physical---that must produce consistent outputs regardless of how inputs are represented.
We do \emph{not} claim to solve the alignment problem, prove metaphysical theses, or provide
guarantees independent of implementation. Rather, we identify \emph{necessary conditions} for
representation-independence and show these conditions have the mathematical structure
of gauge theory. Specifically: (1) evaluation must factor through a declared grounding in
observables; (2) equivalent representations must be identified via canonicalization; (3) evaluation must be invariant under the resulting equivalence relation. When a system satisfies
these conditions and the Lagrangian formulation applies, Noether's theorem implies conservation of a quantity we call the ``alignment current.'' The framework's guarantees are
\emph{conditional} on adequate grounding, correct implementation, and applicability of the continuous symmetry assumptions. We make explicit what the framework does and does not
establish.
\end{abstract}

\vfill

\noindent\rule{\linewidth}{0.4pt}
\small
\noindent\textbf{AI-Assisted Writing Disclosure:} Portions of this manuscript were drafted and refined with AI assistance (Claude, Anthropic). All technical content, mathematical results, and claims were conceived, verified, and approved by the author.

\noindent\textbf{Ethics Statement:} This research involves no human subjects. Empirical validation uses publicly available corpora and synthetic data. No IRB approval was required.

\noindent\textbf{Data Availability:} Code and data are available at \url{https://github.com/ahb-sjsu/erisml-lib}.

\noindent\textbf{Conflict of Interest:} The author declares no conflicts of interest. This research received no external funding.

\noindent\textbf{CRediT Statement:} Andrew H. Bond: Conceptualization, Methodology, Formal Analysis, Investigation, Writing --- Original Draft, Writing --- Review \& Editing, Software.
\noindent\rule{\linewidth}{0.4pt}
\normalsize

\end{titlepage}

\noindent\textbf{Keywords:} representation-independence, gauge theory, grounding axiom, canonicalization, invariant evaluation, fiber bundles, Noether's theorem, AI safety, alignment, ethics

\tableofcontents

\newpage

%=============================================================================
\section{Introduction}
%=============================================================================

\subsection{The Problem}

Consider systems that must evaluate configurations, actions, or states:

\begin{itemize}
    \item An AI system evaluating whether an action is safe
    \item A control system evaluating whether a state satisfies constraints
    \item An ethical reasoning system evaluating whether an action is permissible
\end{itemize}

A common failure mode: the evaluation depends not only on the \emph{situation} but also on how
the situation is \emph{represented}. Changing labels, reordering options, or reformulating descriptions
changes the output, even when the underlying situation is unchanged.

\begin{definition}[Representation-Dependence]\label{def:rep-dep}
An evaluation function $\Sigma : \calX \to V$ is \textbf{representation-dependent} if there exist $x, x' \in \calX$ such that:
\begin{enumerate}
    \item $x$ and $x'$ represent the same underlying situation (by some criterion)
    \item $\Sigma(x) \neq \Sigma(x')$
\end{enumerate}
\end{definition}

\textbf{Goal:} Identify conditions under which evaluation is \emph{representation-independent}. The broader Geometric Ethics framework, which develops these ideas into a full tensorial theory of moral reasoning, is presented in~\cite{bond2026ge}; a reference implementation is available as ErisML/DEME v3.0~\cite{bond2026erisml}.

\subsection{Scope and Limitations}

\textbf{Scope Limitation:} This paper provides \emph{sufficient conditions} for representation-independence
within a specified formal framework. It does \emph{not}:

\begin{itemize}
    \item Prove that these conditions are the \emph{only} way to achieve representation-independence
    \item Guarantee correct behavior if the grounding is inadequate or implementation is faulty
    \item Establish metaphysical claims about the nature of reality, logic, or causation
    \item Solve the ``hard problem'' of choosing which features are morally or physically relevant
\end{itemize}

The framework tells you: \emph{if} you can specify adequate grounding and \emph{if} you implement
correctly, \emph{then} certain guarantees follow.

\subsection{Structure of the Paper}

Section 2 motivates the framework by examining why purely formal approaches face limitations.
Section 3 presents the core definitions and axioms. Section 4 develops the mathematical structure. Section 5 states and proves the main results. Section 6 discusses applications. Section 7
addresses limitations explicitly.

%=============================================================================
\section{Motivation: Limitations of Purely Formal Approaches}
%=============================================================================

We present three observations that motivate the need for grounding. These are not metaphysical
claims but observations about the structure of formal systems.

\subsection{Observation 1: Formal Systems Do Not Determine Their Own Interpretation}

\begin{claim}[Interpretation Underdetermination]\label{claim:interp}
A formal system $F$ consisting of symbols and
inference rules does not, by itself, determine what its symbols refer to or whether its theorems
are true of any particular domain.
\end{claim}

\textbf{Support:} This is a standard result in mathematical logic. Any consistent first-order theory
with infinite models has models of every infinite cardinality (L\"{o}wenheim--Skolem). The same
axioms can be satisfied by vastly different structures. The formal system alone does not pick
out the ``intended'' model.

\textbf{Consequence:} Derivability within a formal system ($F \vdash \phi$) does not establish truth in any
particular domain without an additional specification of how the system is interpreted.

\begin{quote}
\textbf{Scope Limitation:} We are not claiming formal systems are useless or that logic is arbitrary.
We are claiming that formal derivability alone is insufficient to establish claims about specific
domains without interpretation.
\end{quote}

\subsection{Observation 2: G\"{o}del's Theorems Limit Self-Certification}

\begin{theorem}[G\"{o}del's Second Incompleteness Theorem]\label{thm:godel}
If $F$ is a consistent formal system
capable of expressing basic arithmetic, then $F$ cannot prove its own consistency.
\end{theorem}

\textbf{Consequence for our purposes:} A formal system cannot, from within, certify that its
axioms are consistent or that its theorems correspond to truths about an intended domain.
External verification or grounding is required.

\begin{quote}
\textbf{Scope Limitation:} G\"{o}del's theorems apply to sufficiently powerful formal systems (those
capable of expressing arithmetic). They do not apply to all formal systems. We cite them
to illustrate a general point: formal systems have inherent limitations on self-certification.
\end{quote}

\subsection{Observation 3: Multiple Logics Are Internally Coherent}

Classical logic, intuitionistic logic, relevant logic, and various modal logics are all internally
coherent formal systems. They differ in their theorems (e.g., excluded middle is a theorem of
classical but not intuitionistic logic).

\textbf{Consequence:} There is no purely formal argument establishing that one logic is the ``correct'' logic for all domains. The choice of logic is, in part, a choice of tool for a purpose.

\begin{quote}
\textbf{Scope Limitation:} We are not endorsing global logical relativism. Classical logic may well
be correct for most purposes. The point is that the choice cannot be made on purely formal
grounds; it requires some criterion external to the formalisms being compared.
\end{quote}

\subsection{The Upshot}

These observations suggest that trustworthy evaluation requires something beyond formal derivability: a specification of \emph{what the formal system is about} and \emph{how it connects to the domain of
application}. We call this specification \textbf{grounding}.

%=============================================================================
\section{The Formal Framework}
%=============================================================================

We now present the framework precisely. All claims in this section are definitions or axioms,
not empirical claims about the world.

\subsection{Basic Structures}

\begin{definition}[Representation Space]\label{def:rep-space}
A \textbf{representation space} is a set $\calX$ whose elements are
representations (descriptions, encodings, formulations) of situations in some domain of interest.
\end{definition}

\begin{definition}[Value Space]\label{def:val-space}
A \textbf{value space} is a set $V$ (typically $\mathbb{R}$, $\{0,1\}$, or a partial order)
representing possible evaluation outputs.
\end{definition}

\begin{definition}[Evaluation Function]\label{def:eval-func}
An \textbf{evaluation function} is a map $\Sigma : \calX \to V$ assigning
evaluations to representations.
\end{definition}

\subsection{The Grounding Axiom}

The key idea: evaluation should depend on \emph{what is measured}, not on \emph{how it is described}.

\begin{axiom}[Grounding]\label{ax:grounding}
There exists:
\begin{enumerate}[label=(\alph*)]
    \item A \textbf{measurement space} $\calM \subseteq \mathbb{R}^N$ for some $N \in \mathbb{N}$
    \item A \textbf{grounding map} $\Psi : \calX \to \calM$
    \item A \textbf{grounded evaluation} $\Sigmabar : \calM \to V$
\end{enumerate}
such that for all $x \in \calX$:
\[
    \Sigma(x) = \Sigmabar(\Psi(x))
\]
\end{axiom}

\textbf{Interpretation:}
\begin{itemize}
    \item $\calM$ represents measurement outcomes (what sensors report)
    \item $\Psi$ extracts measurements from representations
    \item $\Sigmabar$ operates on measurements, not descriptions
    \item The axiom requires that $\Sigma$ factors through $\Psi$
\end{itemize}

\begin{remark}[What Grounding Does Not Guarantee]\label{rem:grounding-limit}
The Grounding Axiom does not guarantee that $\Psi$ captures all relevant features. If $\Psi$ omits morally or physically relevant aspects
of the situation, evaluation may be ``grounded'' but still inadequate. The adequacy of $\Psi$ is a
domain-specific question outside the scope of the formal framework.
\end{remark}

\subsection{The Equivalence Relation}

Grounding induces an equivalence relation on representations.

\begin{definition}[Grounding Equivalence]\label{def:grounding-equiv}
Two representations $x, x' \in \calX$ are \textbf{grounding-equivalent}, written $x \sim_\Psi x'$, if and only if $\Psi(x) = \Psi(x')$.
\end{definition}

\begin{proposition}[Invariance Under Equivalence]\label{prop:inv-equiv}
If Axiom~\ref{ax:grounding} holds, then $\Sigma$ is constant on
equivalence classes:
\[
    x \sim_\Psi x' \implies \Sigma(x) = \Sigma(x')
\]
\end{proposition}

\begin{proof}
$x \sim_\Psi x'$ means $\Psi(x) = \Psi(x')$. By Axiom~\ref{ax:grounding}, $\Sigma(x) = \Sigmabar(\Psi(x)) = \Sigmabar(\Psi(x')) = \Sigma(x')$.
\end{proof}

\subsection{The Canonicalization Axiom}

A canonicalizer selects a unique representative from each equivalence class.

\begin{axiom}[Canonicalization]\label{ax:canon}
There exists a \textbf{canonicalizer} $\kappa : \calX \to \calX$ satisfying:
\begin{enumerate}[label=(\alph*)]
    \item \textbf{Idempotence:} $\kappa(\kappa(x)) = \kappa(x)$ for all $x \in \calX$
    \item \textbf{Equivalence-class collapse:} $x \sim_\Psi x' \implies \kappa(x) = \kappa(x')$
    \item \textbf{Grounding preservation:} $\Psi(\kappa(x)) = \Psi(x)$ for all $x \in \calX$
\end{enumerate}
\end{axiom}

\begin{proposition}[Canonicalization Yields Representative]\label{prop:canon-rep}
Under Axiom~\ref{ax:canon}, $\kappa(x)$ is a well-defined representative of the equivalence class $[x]_\Psi$.
\end{proposition}

\begin{proof}
By (b), $\kappa$ maps all elements of $[x]_\Psi$ to the same element. By (c), this element is in $[x]_\Psi$.
By (a), it is a fixed point of $\kappa$.
\end{proof}

\subsection{The Transformation Group}

\begin{definition}[Structure-Preserving Transformations]\label{def:struct-group}
The \textbf{structure group} $\calG$ is the set of
bijections $g : \calX \to \calX$ that preserve grounding:
\[
    \calG := \{g : \calX \to \calX \mid g \text{ is a bijection and } \Psi(g(x)) = \Psi(x) \text{ for all } x\}
\]
\end{definition}

\begin{proposition}[$\calG$ is a Group]\label{prop:group}
$\calG$ forms a group under composition.
\end{proposition}

\begin{proof}
Identity: $\mathrm{id} \in \calG$ (trivially preserves $\Psi$). Closure: if $g, h \in \calG$, then $\Psi(gh(x)) = \Psi(h(x)) =
\Psi(x)$. Inverse: if $g \in \calG$ and $y = g(x)$, then $\Psi(g^{-1}(y)) = \Psi(x) = \Psi(g(x)) = \Psi(y)$.
\end{proof}

\subsection{The Invariance Theorem}

\begin{theorem}[Invariance]\label{thm:invariance}
If Axiom~\ref{ax:grounding} holds, then $\Sigma$ is $\calG$-invariant:
\[
    \Sigma(g(x)) = \Sigma(x) \quad \forall g \in \calG,\; x \in \calX
\]
\end{theorem}

\begin{proof}
Let $g \in \calG$. By definition, $\Psi(g(x)) = \Psi(x)$. By Axiom~\ref{ax:grounding}, $\Sigma(g(x)) = \Sigmabar(\Psi(g(x))) =
\Sigmabar(\Psi(x)) = \Sigma(x)$.
\end{proof}

\begin{remark}[Converse]\label{rem:converse}
The converse is also true under mild conditions: if $\Sigma$ is $\calG$-invariant
and $\calG$ acts transitively on $\Psi$-equivalence classes, then $\Sigma$ factors through $\Psi$.
\end{remark}

%=============================================================================
\section{The Geometric Structure}
%=============================================================================

We now show that the framework has the structure of a principal fiber bundle. This section
assumes familiarity with differential geometry; readers primarily interested in applications may
skip to Section~5.

\subsection{Bundle Structure}

\begin{definition}[Principal Bundle Conditions]\label{def:bundle-cond}
The quadruple $(\calX, \calM, \Psi, \calG)$ satisfies \textbf{principal
bundle conditions} if:
\end{definition}

\begin{enumerate}[label=(PB\arabic*)]
    \item $\calG$ acts freely on $\calX$: $g(x) = x \implies g = \mathrm{id}$
    \item $\calG$ acts transitively on fibers: $\Psi(x) = \Psi(x') \implies \exists g \in \calG : g(x) = x'$
    \item $\Psi : \calX \to \calM$ is surjective
    \item (Smoothness, if applicable): $\calX$, $\calM$ are smooth manifolds; $\calG$ is a Lie group; $\Psi$ is a smooth submersion
\end{enumerate}

\begin{theorem}[Bundle Structure]\label{thm:bundle}
Under conditions (PB1)--(PB3), $\Psi : \calX \to \calM$ is a principal
$\calG$-bundle in the topological sense. Under (PB1)--(PB4), it is a smooth principal bundle.
\end{theorem}

\begin{proof}
(PB1) gives freeness. (PB2) implies orbits equal fibers. (PB3) gives surjectivity. The
quotient $\calX/\calG \cong \calM$ follows. For the smooth case, (PB4) plus standard results on Lie group
actions give local triviality.
\end{proof}

\begin{quote}
\textbf{Scope Limitation:} The bundle structure requires (PB1)--(PB2), which are substantive
assumptions. Not all grounding setups satisfy them. When they fail, the geometric picture
is only partially applicable.
\end{quote}

\subsection{Connection and Curvature}

\begin{definition}[Connection via Section]\label{def:connection}
A global section $\sigma : \calM \to \calX$ satisfying $\Psi \circ \sigma = \mathrm{id}_\calM$
determines a connection on the bundle.
\end{definition}

\begin{theorem}[Canonicalization Gives Section]\label{thm:canon-section}
Under Axioms~\ref{ax:grounding} and~\ref{ax:canon}, if the bundle conditions
(PB1)--(PB2) hold, then $\kappa$ induces a global section $\sigma : \calM \to \calX$ defined by:
\[
    \sigma(m) := \kappa(x) \quad \text{for any } x \in \Psi^{-1}(m)
\]
\end{theorem}

\begin{proof}
By Axiom~\ref{ax:canon}(b), $\kappa(x) = \kappa(x')$ whenever $\Psi(x) = \Psi(x')$, so $\sigma$ is well-defined. By Axiom~\ref{ax:canon}(c), $\Psi(\sigma(m)) = \Psi(\kappa(x)) = \Psi(x) = m$, so $\Psi \circ \sigma = \mathrm{id}_\calM$.
\end{proof}

\begin{theorem}[Flatness]\label{thm:flatness}
The connection determined by a global section has zero curvature.
\end{theorem}

\begin{proof}
A global section trivializes the bundle: $\calX \cong \calM \times \calG$ via $(m, g) \mapsto \sigma(m) \cdot g$. In a trivial
bundle, the canonical flat connection has curvature $\Omega = d\omega + \frac{1}{2}[\omega, \omega] = 0$ (Maurer--Cartan
equation on the $\calG$ factor).
\end{proof}

\begin{remark}[Interpretation of Flatness]\label{rem:flatness}
Flatness means parallel transport is path-independent:
representations can be ``moved'' between fibers without holonomy. Physically, this reflects that
representational choices have no intrinsic structure---they are ``pure gauge.''
\end{remark}

%=============================================================================
\section{Conservation Results}
%=============================================================================

This section shows that, under additional assumptions, invariance implies conservation. The
results require a Lagrangian formulation, which may not apply in all contexts.

\subsection{The Lagrangian Setting}

\begin{definition}[Decision Lagrangian]\label{def:lagrangian}
Suppose $\calX$ and $\calM$ are smooth manifolds. A \textbf{decision
Lagrangian} is a function $L : T\calX \to \mathbb{R}$ of the form:
\begin{equation}\label{eq:lagrangian}
    L(x, \dot{x}) = K(\dot{x}) - U(\Psi(x)) - \frac{\lambda}{2}\|\omega(\dot{x})\|^2_\fg
\end{equation}
where:
\begin{itemize}
    \item $K$ is a kinetic term (cost of deliberation)
    \item $U : \calM \to \mathbb{R}$ is a potential (negative of grounded evaluation)
    \item $\omega$ is the connection form
    \item $\lambda > 0$ is a coupling constant
    \item $\|\cdot\|_\fg$ is an -invariant norm on the Lie algebra $\fg$
\end{itemize}
\end{definition}

\begin{quote}
\textbf{Scope Limitation:} The Lagrangian formulation assumes continuous dynamics on smooth
manifolds. It applies to systems whose evolution can be described by differential equations.
Discrete systems, non-smooth dynamics, and systems without natural ``kinetic'' interpretations require separate analysis.
\end{quote}

\subsection{Noether's Theorem}

\begin{theorem}[Noether]\label{thm:noether}
Let $L : TQ \to \mathbb{R}$ be a Lagrangian on configuration space $Q$. If a
one-parameter group $\{\phi_s\}_{s \in \mathbb{R}}$ satisfies $L \circ d\phi_s = L$ (symmetry), then the quantity
\[
    Q_\xi := \left\langle \frac{\partial L}{\partial \dot{q}}, \xi \right\rangle
\]
is constant along solutions of the Euler--Lagrange equations, where $\xi = \frac{d}{ds}\big|_{s=0} \phi_s(q)$ is the infinitesimal generator.
\end{theorem}

\begin{proof}
Standard. See Arnold, \emph{Mathematical Methods of Classical Mechanics}, Chapter 4.
\end{proof}

\subsection{Conservation of Alignment}

\begin{theorem}[Alignment Current]\label{thm:alignment}
For the Lagrangian~\eqref{eq:lagrangian}, the gauge symmetry $x \mapsto x \cdot g$
implies a conserved current:
\[
    J := \lambda\omega(\dot{x}) \in \fg
\]
satisfying $\frac{d}{dt}J = 0$ along Euler--Lagrange trajectories.
\end{theorem}

\begin{proof}
The Lagrangian is $\calG$-invariant by construction: $K$ depends only on velocity magnitude;
$U$ depends only on $\Psi(x)$ which is $\calG$-invariant; $\|\omega(\dot{x})\|$ is -invariant. Apply Theorem~\ref{thm:noether} to the
one-parameter subgroups of $\calG$.
\end{proof}

\begin{corollary}[Horizontal Trajectories are Preserved]\label{cor:horizontal}
If $\omega(\dot{x}(0)) = 0$ (trajectory starts horizontal), then $\omega(\dot{x}(t)) = 0$ for all $t$.
\end{corollary}

\begin{proof}
$J(t) = \lambda\omega(\dot{x}(t)) = J(0) = 0$ by conservation.
\end{proof}

\begin{remark}[Interpretation]\label{rem:horizontal}
``Horizontal'' means the trajectory has no component in the
fiber direction---it moves only in the base (physical) directions. Corollary~\ref{cor:horizontal} says: \emph{a system
that starts responding only to physical states, not to representational artifacts, continues to do
so}.
\end{remark}

\begin{quote}
\textbf{Scope Limitation:} Conservation requires: (1) the Lagrangian formulation applies, (2) the
symmetry is continuous (Lie group), (3) dynamics follow Euler--Lagrange equations. For
discrete symmetry groups (e.g., permutations) or non-Lagrangian dynamics, conservation in
this form does not apply. The invariance results (Section 4) still hold; only the dynamical
conservation fails.
\end{quote}

%=============================================================================
\section{Applications}
%=============================================================================

\subsection{Application to Ethics: The Bond Invariance Principle}

\begin{definition}[Bond]\label{def:bond}
A \textbf{bond} is a morally relevant relationship between entities, represented
as a tuple $b = (a, p, r, c)$ where $a$ is an agent, $p$ is a patient, $r$ is a relation type, and $c$ is context.
\end{definition}

\begin{definition}[Bond Structure]\label{def:bond-structure}
The \textbf{bond structure} $B(x)$ of a representation $x$ is the set
of all bonds encoded in $x$.
\end{definition}

\textbf{Instantiation of the framework:}
\begin{itemize}
    \item $\Psi(x) := B(x)$ (grounding in bond structure)
    \item $\calG :=$ bond-preserving transformations (relabeling, reordering that preserve $B$)
    \item $\Sigma :=$ moral evaluation function
\end{itemize}

\begin{corollary}[Bond Invariance Principle]\label{cor:bip}
If moral evaluation satisfies Axiom~\ref{ax:grounding} with $\Psi = B$,
then:
\[
    B(x) = B(x') \implies \Sigma(x) = \Sigma(x')
\]
Moral judgments are invariant under bond-preserving transformations.
\end{corollary}

\begin{quote}
\textbf{Scope Limitation:} This does not tell us \emph{which} relationships are morally relevant (the
content of $B$). It says: \emph{whatever you count as morally relevant, evaluation should depend
only on that, not on representational choices}. The choice of $B$ is a substantive ethical
question outside the formal framework.
\end{quote}

\subsection{Application to Control Systems}

\textbf{Instantiation for plasma control:}
\begin{itemize}
    \item $\Psi(x) := (\bar{n}, I_p, \beta_N, \ldots)$ (measured plasma parameters)
    \item $\calG :=$ unit changes, coordinate transformations
    \item $\Sigma :=$ control objective / safety score
\end{itemize}

\begin{corollary}[Unit Invariance]\label{cor:unit-inv}
If a controller satisfies Axiom~\ref{ax:grounding}, its outputs are independent
of unit conventions.
\end{corollary}

\begin{corollary}[Dimensionless Transfer]\label{cor:dimless}
If $\Psi$ consists of dimensionless quantities, controllers
transfer across systems with different physical scales.
\end{corollary}

\subsection{Application to AI Safety}

\textbf{Instantiation for AI containment:}
\begin{itemize}
    \item $\Psi(x) :=$ sensor readings, physical observables
    \item $\calG :=$ representational manipulations (relabeling, redescription)
    \item $\Sigma :=$ safety/alignment score
\end{itemize}

\begin{theorem}[Representation Attack Immunity]\label{thm:attack-immune}
If an AI system's evaluation satisfies Axiom~\ref{ax:grounding}, then for any $\Psi$-preserving manipulation $g \in \calG$:
\[
    \Sigma(g(x)) = \Sigma(x)
\]
The manipulation does not change the evaluation.
\end{theorem}

\begin{proof}
Immediate from Theorem~\ref{thm:invariance}.
\end{proof}

\begin{quote}
\textbf{Scope Limitation:} Theorem~\ref{thm:attack-immune} blocks \emph{representational} attacks (relabeling, redescription).
It does \emph{not} block:
\begin{itemize}
    \item Physical attacks (sensor spoofing)
    \item Attacks that change $\Psi$-values (actual physical manipulation)
    \item Attacks on the implementation of $\kappa$, $\Psi$, or $\Sigma$
    \item Inadequacy of $\Psi$ (missing relevant features)
\end{itemize}
The theorem's scope is precisely: \emph{given correct implementation and adequate grounding,
representational evasion fails}.
\end{quote}

%=============================================================================
\section{Limitations and Non-Claims}
%=============================================================================

We explicitly state what the framework does \emph{not} establish.

\subsection{The Framework Does Not Solve the Grounding Problem}

\begin{claim}[Grounding is Assumed, Not Derived]\label{claim:grounding-assumed}
The framework assumes a grounding map $\Psi$
is given. It provides no procedure for:
\begin{itemize}
    \item Determining which features are morally/physically relevant
    \item Verifying that $\Psi$ is adequate for the domain
    \item Discovering that $\Psi$ is missing important features
\end{itemize}
\end{claim}

\textbf{Consequence:} A system can satisfy all formal requirements of the framework while being
inadequately grounded. The framework makes explicit \emph{what must be specified} ($\Psi$) and \emph{what
properties follow if it is specified correctly}, but the correctness of $\Psi$ itself is a domain-specific
judgment.

\subsection{The Framework Does Not Guarantee Implementation Correctness}

\begin{claim}[Implementation is Assumed, Not Verified]\label{claim:impl}
The framework's guarantees are conditional on:
\begin{itemize}
    \item Correct implementation of $\kappa$, $\Psi$, and $\Sigma$
    \item Correct implementation of the trusted computing base
    \item Physical security of sensors and actuators
\end{itemize}
\end{claim}

\textbf{Consequence:} Bugs, hardware failures, or security breaches can violate the guarantees.
The framework is a design specification, not a correctness proof for any particular implementation.

\subsection{Conservation Requires Continuous Symmetry and Lagrangian Dynamics}

\begin{claim}[Conservation Has Preconditions]\label{claim:conservation}
The conservation results (Theorem~\ref{thm:alignment}, Corollary~\ref{cor:horizontal}) require:
\begin{itemize}
    \item $\calG$ is a Lie group (continuous symmetry)
    \item Dynamics are governed by a Lagrangian
    \item The system evolves according to Euler--Lagrange equations
\end{itemize}
\end{claim}

For discrete $\calG$ (e.g., permutation groups) or non-Lagrangian dynamics, conservation in this form
does not hold.

\textbf{Consequence:} The invariance results (Theorem~\ref{thm:invariance}) are more general than the conservation results. Invariance holds whenever Axiom~\ref{ax:grounding} holds; conservation requires additional
structure.

\subsection{The Framework Does Not Establish Metaphysical Claims}

\begin{claim}[No Metaphysical Commitments]\label{claim:metaphysics}
The framework does not establish:
\begin{itemize}
    \item That reality has any particular structure
    \item That logic, causation, or modality work in any particular way
    \item That grounding in observables is the only way to achieve trustworthiness
\end{itemize}
\end{claim}

The framework is a \emph{conditional}: if you want representation-independence, and if you can
specify adequate grounding, then these mathematical structures provide it. The framework is
compatible with many metaphysical views.

\subsection{Summary of Claims}

\begin{table}[h]
\centering
\begin{tabular}{lll}
\toprule
\textbf{Claim} & \textbf{Established?} & \textbf{Requires} \\
\midrule
Grounding $\Rightarrow$ Invariance & Yes & Axiom~\ref{ax:grounding} \\
Bundle structure exists & Yes & (PB1)--(PB3) \\
Canonicalization gives section & Yes & Axioms~\ref{ax:grounding},~\ref{ax:canon}, (PB1)--(PB2) \\
Connection is flat & Yes & Global section exists \\
Alignment current is conserved & Yes & Lagrangian, Lie group $\calG$, E-L dynamics \\
$\Psi$ is adequate & \textbf{No} & Domain expertise \\
Implementation is correct & \textbf{No} & Verification \\
All attacks are blocked & \textbf{No} & Only representational attacks \\
\bottomrule
\end{tabular}
\end{table}

%=============================================================================
\section{Conclusion}
%=============================================================================

\subsection{What We Have Established}

\begin{enumerate}
    \item \textbf{Definition:} Representation-independence means evaluation factors through a grounding
    map.
    \item \textbf{Sufficient conditions:} Axioms~\ref{ax:grounding} and~\ref{ax:canon} suffice for representation-independence.
    \item \textbf{Mathematical structure:} Under principal bundle conditions, the framework has gauge-theoretic structure with flat connection.
    \item \textbf{Conservation:} Under Lagrangian dynamics with continuous symmetry, alignment is
    conserved.
    \item \textbf{Applications:} The framework instantiates in ethics (BIP), control (unit invariance), and
    AI safety (representation attack immunity).
\end{enumerate}

\subsection{What Remains Open}

\begin{enumerate}
    \item \textbf{Grounding selection:} How to choose $\Psi$ for a given domain.
    \item \textbf{Adequacy verification:} How to determine if $\Psi$ is adequate.
    \item \textbf{Implementation:} How to correctly implement $\kappa$, $\Psi$, $\Sigma$.
    \item \textbf{Discrete symmetries:} Conservation-like results for non-Lie groups.
    \item \textbf{Non-Lagrangian systems:} Dynamical guarantees without variational structure.
\end{enumerate}

\subsection{The Conservative Claim}

We make a conservative claim:

\begin{quote}
\textbf{If} you can specify a grounding map $\Psi$ that captures the relevant features of your domain, \textbf{and if} you correctly implement canonicalization and grounded evaluation, \textbf{then} your
evaluation will be representation-independent in the precise sense that $\calG$-transformations
do not change outputs.

Under additional assumptions (Lagrangian dynamics, continuous symmetry), alignment
is conserved: systems that start responding only to grounded features continue to do so.

The framework does not tell you what to ground in, does not verify your implementation,
and does not block non-representational attacks. It provides structure, not magic.
\end{quote}

%=============================================================================
\appendix
%=============================================================================

\section{Proof of Theorem~\ref{thm:bundle}}\label{app:bundle-proof}

\begin{proof}
We verify the principal bundle axioms.

\textbf{Free action:} (PB1) states $g(x) = x \implies g = \mathrm{id}$, which is the definition of free action.

\textbf{Transitive action on fibers:} (PB2) states that for any $x, x'$ with $\Psi(x) = \Psi(x')$, there
exists $g$ with $g(x) = x'$. This means $\calG$ acts transitively on each fiber $\Psi^{-1}(m)$.

\textbf{Orbits equal fibers:} If $x' = g(x)$ for some $g \in \calG$, then $\Psi(x') = \Psi(g(x)) = \Psi(x)$ (since $g$
preserves $\Psi$). So orbits are contained in fibers. By (PB2), fibers are contained in orbits. Hence
orbits equal fibers.

\textbf{Base space:} $\calM = \calX/\calG$ follows from orbits equaling fibers. (PB3) ensures $\Psi$ is surjective.

\textbf{Local triviality (smooth case):} Under (PB4), $\Psi$ is a smooth submersion with $\calG$ a Lie
group acting freely and properly. Standard results (e.g., Lee, \emph{Introduction to Smooth Manifolds},
Theorem 21.10) give local triviality.
\end{proof}

\section{Relationship to Physics Literature}\label{app:physics}

The gauge-theoretic structure we describe is standard in differential geometry and mathematical
physics:

\begin{itemize}
    \item \textbf{Principal bundles:} Kobayashi \& Nomizu, \emph{Foundations of Differential Geometry}
    \item \textbf{Connections and curvature:} Nakahara, \emph{Geometry, Topology and Physics}
    \item \textbf{Noether's theorem:} Arnold, \emph{Mathematical Methods of Classical Mechanics}
    \item \textbf{Gauge theory in physics:} Bleecker, \emph{Gauge Theory and Variational Principles}
\end{itemize}

Our contribution is not to the mathematics of gauge theory but to its application as a
framework for representation-independent evaluation.

\section{Glossary}\label{app:glossary}

\begin{table}[h]
\centering
\begin{tabularx}{\textwidth}{lX}
\toprule
\textbf{Term} & \textbf{Definition} \\
\midrule
Grounding map $\Psi$ & Function extracting measurements from representations \\
Canonicalizer $\kappa$ & Function selecting unique representative per equivalence class \\
Structure group $\calG$ & Group of $\Psi$-preserving transformations \\
Grounding equivalence $\sim_\Psi$ & $x \sim_\Psi x'$ iff $\Psi(x) = \Psi(x')$ \\
Principal bundle & Fiber bundle where fibers are group orbits \\
Connection & Specification of ``horizontal'' directions in bundle \\
Flat connection & Connection with zero curvature \\
Noether current & Conserved quantity from continuous symmetry \\
BIP & Bond Invariance Principle: moral invariance under bond-preserving transformations \\
\bottomrule
\end{tabularx}
\end{table}

%=============================================================================
% Bibliography
%=============================================================================

\begin{thebibliography}{9}

\bibitem{arnold}
V.~I. Arnold.
\newblock \emph{Mathematical Methods of Classical Mechanics}.
\newblock Springer, 2nd edition, 1989.

\bibitem{bleecker}
D.~Bleecker.
\newblock \emph{Gauge Theory and Variational Principles}.
\newblock Dover, 2005.

\bibitem{kobayashi}
S.~Kobayashi and K.~Nomizu.
\newblock \emph{Foundations of Differential Geometry}, Volumes I and II.
\newblock Wiley, 1963/1969.

\bibitem{lee}
J.~M. Lee.
\newblock \emph{Introduction to Smooth Manifolds}.
\newblock Springer, 2nd edition, 2012.

\bibitem{nakahara}
M.~Nakahara.
\newblock \emph{Geometry, Topology and Physics}.
\newblock CRC Press, 2nd edition, 2003.

\bibitem{bond2026ge}
A.~H.~Bond, \emph{Geometric Ethics: The Mathematical Structure of Moral Reasoning}, v1.14, Department of Computer Engineering, San Jos\'e State University, Spring 2026.

\bibitem{bond2026erisml}
A.~H.~Bond, ``ErisML/DEME v3.0: A Library for Governed, Foundation-Model-Enabled Agents with Tensorial Ethics,'' \url{https://github.com/ahb-sjsu/erisml-lib}, 2026.

\end{thebibliography}

\end{document}
